\documentclass[
  doc,
  floatsintext,
  longtable,
  a4paper,
  nolmodern,
  notxfonts,
  notimes,
  12pt,
  colorlinks=true,linkcolor=blue,citecolor=blue,urlcolor=blue]{apa7}

\usepackage{amsmath}
\usepackage{amssymb}

\geometry{inner=1in, outer=1in}
\fancyhfoffset[LE,RO]{0cm}


\usepackage[bidi=default]{babel}
\babelprovide[main,import]{spanish}


% get rid of language-specific shorthands (see #6817):
\let\LanguageShortHands\languageshorthands
\def\languageshorthands#1{}

\RequirePackage{longtable}
\RequirePackage{threeparttablex}

\makeatletter
\renewcommand{\paragraph}{\@startsection{paragraph}{4}{\parindent}%
	{0\baselineskip \@plus 0.2ex \@minus 0.2ex}%
	{-.5em}%
	{\normalfont\normalsize\bfseries\typesectitle}}

\renewcommand{\subparagraph}[1]{\@startsection{subparagraph}{5}{0.5em}%
	{0\baselineskip \@plus 0.2ex \@minus 0.2ex}%
	{-\z@\relax}%
	{\normalfont\normalsize\bfseries\itshape\hspace{\parindent}{#1}\textit{\addperi}}{\relax}}
\makeatother




\usepackage{longtable, booktabs, multirow, multicol, colortbl, hhline, caption, array, float, xpatch}
\usepackage{subcaption}


\renewcommand\thesubfigure{\Alph{subfigure}}
\setcounter{topnumber}{2}
\setcounter{bottomnumber}{2}
\setcounter{totalnumber}{4}
\renewcommand{\topfraction}{0.85}
\renewcommand{\bottomfraction}{0.85}
\renewcommand{\textfraction}{0.15}
\renewcommand{\floatpagefraction}{0.7}

\usepackage{tcolorbox}
\tcbuselibrary{listings,theorems, breakable, skins}
\usepackage{fontawesome5}

\definecolor{quarto-callout-color}{HTML}{909090}
\definecolor{quarto-callout-note-color}{HTML}{0758E5}
\definecolor{quarto-callout-important-color}{HTML}{CC1914}
\definecolor{quarto-callout-warning-color}{HTML}{EB9113}
\definecolor{quarto-callout-tip-color}{HTML}{00A047}
\definecolor{quarto-callout-caution-color}{HTML}{FC5300}
\definecolor{quarto-callout-color-frame}{HTML}{ACACAC}
\definecolor{quarto-callout-note-color-frame}{HTML}{4582EC}
\definecolor{quarto-callout-important-color-frame}{HTML}{D9534F}
\definecolor{quarto-callout-warning-color-frame}{HTML}{F0AD4E}
\definecolor{quarto-callout-tip-color-frame}{HTML}{02B875}
\definecolor{quarto-callout-caution-color-frame}{HTML}{FD7E14}

%\newlength\Oldarrayrulewidth
%\newlength\Oldtabcolsep


\usepackage{hyperref}



\usepackage{color}
\usepackage{fancyvrb}
\newcommand{\VerbBar}{|}
\newcommand{\VERB}{\Verb[commandchars=\\\{\}]}
\DefineVerbatimEnvironment{Highlighting}{Verbatim}{commandchars=\\\{\}}
% Add ',fontsize=\small' for more characters per line
\usepackage{framed}
\definecolor{shadecolor}{RGB}{241,243,245}
\newenvironment{Shaded}{\begin{snugshade}}{\end{snugshade}}
\newcommand{\AlertTok}[1]{\textcolor[rgb]{0.68,0.00,0.00}{#1}}
\newcommand{\AnnotationTok}[1]{\textcolor[rgb]{0.37,0.37,0.37}{#1}}
\newcommand{\AttributeTok}[1]{\textcolor[rgb]{0.40,0.45,0.13}{#1}}
\newcommand{\BaseNTok}[1]{\textcolor[rgb]{0.68,0.00,0.00}{#1}}
\newcommand{\BuiltInTok}[1]{\textcolor[rgb]{0.00,0.23,0.31}{#1}}
\newcommand{\CharTok}[1]{\textcolor[rgb]{0.13,0.47,0.30}{#1}}
\newcommand{\CommentTok}[1]{\textcolor[rgb]{0.37,0.37,0.37}{#1}}
\newcommand{\CommentVarTok}[1]{\textcolor[rgb]{0.37,0.37,0.37}{\textit{#1}}}
\newcommand{\ConstantTok}[1]{\textcolor[rgb]{0.56,0.35,0.01}{#1}}
\newcommand{\ControlFlowTok}[1]{\textcolor[rgb]{0.00,0.23,0.31}{\textbf{#1}}}
\newcommand{\DataTypeTok}[1]{\textcolor[rgb]{0.68,0.00,0.00}{#1}}
\newcommand{\DecValTok}[1]{\textcolor[rgb]{0.68,0.00,0.00}{#1}}
\newcommand{\DocumentationTok}[1]{\textcolor[rgb]{0.37,0.37,0.37}{\textit{#1}}}
\newcommand{\ErrorTok}[1]{\textcolor[rgb]{0.68,0.00,0.00}{#1}}
\newcommand{\ExtensionTok}[1]{\textcolor[rgb]{0.00,0.23,0.31}{#1}}
\newcommand{\FloatTok}[1]{\textcolor[rgb]{0.68,0.00,0.00}{#1}}
\newcommand{\FunctionTok}[1]{\textcolor[rgb]{0.28,0.35,0.67}{#1}}
\newcommand{\ImportTok}[1]{\textcolor[rgb]{0.00,0.46,0.62}{#1}}
\newcommand{\InformationTok}[1]{\textcolor[rgb]{0.37,0.37,0.37}{#1}}
\newcommand{\KeywordTok}[1]{\textcolor[rgb]{0.00,0.23,0.31}{\textbf{#1}}}
\newcommand{\NormalTok}[1]{\textcolor[rgb]{0.00,0.23,0.31}{#1}}
\newcommand{\OperatorTok}[1]{\textcolor[rgb]{0.37,0.37,0.37}{#1}}
\newcommand{\OtherTok}[1]{\textcolor[rgb]{0.00,0.23,0.31}{#1}}
\newcommand{\PreprocessorTok}[1]{\textcolor[rgb]{0.68,0.00,0.00}{#1}}
\newcommand{\RegionMarkerTok}[1]{\textcolor[rgb]{0.00,0.23,0.31}{#1}}
\newcommand{\SpecialCharTok}[1]{\textcolor[rgb]{0.37,0.37,0.37}{#1}}
\newcommand{\SpecialStringTok}[1]{\textcolor[rgb]{0.13,0.47,0.30}{#1}}
\newcommand{\StringTok}[1]{\textcolor[rgb]{0.13,0.47,0.30}{#1}}
\newcommand{\VariableTok}[1]{\textcolor[rgb]{0.07,0.07,0.07}{#1}}
\newcommand{\VerbatimStringTok}[1]{\textcolor[rgb]{0.13,0.47,0.30}{#1}}
\newcommand{\WarningTok}[1]{\textcolor[rgb]{0.37,0.37,0.37}{\textit{#1}}}

\providecommand{\tightlist}{%
  \setlength{\itemsep}{0pt}\setlength{\parskip}{0pt}}
\usepackage{longtable,booktabs,array}
\usepackage{calc} % for calculating minipage widths
% Correct order of tables after \paragraph or \subparagraph
\usepackage{etoolbox}
\makeatletter
\patchcmd\longtable{\par}{\if@noskipsec\mbox{}\fi\par}{}{}
\makeatother
% Allow footnotes in longtable head/foot
\IfFileExists{footnotehyper.sty}{\usepackage{footnotehyper}}{\usepackage{footnote}}
\makesavenoteenv{longtable}

\usepackage{graphicx}
\makeatletter
\newsavebox\pandoc@box
\newcommand*\pandocbounded[1]{% scales image to fit in text height/width
  \sbox\pandoc@box{#1}%
  \Gscale@div\@tempa{\textheight}{\dimexpr\ht\pandoc@box+\dp\pandoc@box\relax}%
  \Gscale@div\@tempb{\linewidth}{\wd\pandoc@box}%
  \ifdim\@tempb\p@<\@tempa\p@\let\@tempa\@tempb\fi% select the smaller of both
  \ifdim\@tempa\p@<\p@\scalebox{\@tempa}{\usebox\pandoc@box}%
  \else\usebox{\pandoc@box}%
  \fi%
}
% Set default figure placement to htbp
\def\fps@figure{htbp}
\makeatother







\usepackage{newtx}

\defaultfontfeatures{Scale=MatchLowercase}
\defaultfontfeatures[\rmfamily]{Ligatures=TeX,Scale=1}





\title{Instalación y Configuración de Miniconda}


\shorttitle{Instalación y Configuración de Miniconda}


\usepackage{etoolbox}



\ccoppy{\textcopyright~2020}



\author{Edison Achalma}



\affiliation{
{Escuela Profesional de Economía, Universidad Nacional de San Cristóbal
de Huamanga}}




\leftheader{Achalma}

\date{2020-06-19}


\abstract{Este abstract será actualizado una vez que se complete el
contenido final del artículo. }

\keywords{keyword1, keyword2}

\authornote{\par{\addORCIDlink{Edison Achalma}{0000-0001-6996-3364}} 
\par{ }
\par{   El autor no tiene conflictos de interés que revelar.    Los
roles de autor se clasificaron utilizando la taxonomía de roles de
colaborador (CRediT; https://credit.niso.org/) de la siguiente
manera:  Edison Achalma:   conceptualización, redacción}
\par{La correspondencia relativa a este artículo debe dirigirse a Edison
Achalma, Email: \href{mailto:elmer.achalma.09@unsch.edu.pe}{elmer.achalma.09@unsch.edu.pe}}
}

\makeatletter
\let\endoldlt\endlongtable
\def\endlongtable{
\hline
\endoldlt
}
\makeatother

\urlstyle{same}



\makeatletter
\@ifpackageloaded{caption}{}{\usepackage{caption}}
\AtBeginDocument{%
\ifdefined\contentsname
  \renewcommand*\contentsname{Tabla de contenidos}
\else
  \newcommand\contentsname{Tabla de contenidos}
\fi
\ifdefined\listfigurename
  \renewcommand*\listfigurename{Lista de Figuras}
\else
  \newcommand\listfigurename{Lista de Figuras}
\fi
\ifdefined\listtablename
  \renewcommand*\listtablename{Lista de Tablas}
\else
  \newcommand\listtablename{Lista de Tablas}
\fi
\ifdefined\figurename
  \renewcommand*\figurename{Figura}
\else
  \newcommand\figurename{Figura}
\fi
\ifdefined\tablename
  \renewcommand*\tablename{Tabla}
\else
  \newcommand\tablename{Tabla}
\fi
}
\@ifpackageloaded{float}{}{\usepackage{float}}
\floatstyle{ruled}
\@ifundefined{c@chapter}{\newfloat{codelisting}{h}{lop}}{\newfloat{codelisting}{h}{lop}[chapter]}
\floatname{codelisting}{Listing}
\newcommand*\listoflistings{\listof{codelisting}{Lista de Listings}}
\makeatother
\makeatletter
\makeatother
\makeatletter
\@ifpackageloaded{caption}{}{\usepackage{caption}}
\@ifpackageloaded{subcaption}{}{\usepackage{subcaption}}
\makeatother
\makeatletter
\@ifpackageloaded{fontawesome5}{}{\usepackage{fontawesome5}}
\makeatother

% From https://tex.stackexchange.com/a/645996/211326
%%% apa7 doesn't want to add appendix section titles in the toc
%%% let's make it do it
\makeatletter
\xpatchcmd{\appendix}
  {\par}
  {\addcontentsline{toc}{section}{\@currentlabelname}\par}
  {}{}
\makeatother

%% Disable longtable counter
%% https://tex.stackexchange.com/a/248395/211326

\usepackage{etoolbox}

\makeatletter
\patchcmd{\LT@caption}
  {\bgroup}
  {\bgroup\global\LTpatch@captiontrue}
  {}{}
\patchcmd{\longtable}
  {\par}
  {\par\global\LTpatch@captionfalse}
  {}{}
\apptocmd{\endlongtable}
  {\ifLTpatch@caption\else\addtocounter{table}{-1}\fi}
  {}{}
\newif\ifLTpatch@caption
\makeatother

\begin{document}

\maketitle


\hypertarget{toc}{}
\tableofcontents
\newpage
\section[Introduction]{Instalación y Configuración de Miniconda}

\setcounter{secnumdepth}{3}

\setlength\LTleft{0pt}




\section{¿Qué es Anaconda?}\label{quuxe9-es-anaconda}

Anaconda es una plataforma de gestión de paquetes y entornos virtuales
diseñada para facilitar el desarrollo en Python y otros lenguajes de
programación populares. Imagina que tienes una caja de herramientas
repleta de todo lo que necesitas para construir aplicaciones, analizar
datos o desarrollar proyectos de aprendizaje automático. Esa es
Anaconda: una caja de herramientas poderosa y completa que te brinda
todo lo que necesitas para trabajar con Python de manera eficiente.

Esta guía está diseñada específicamente para mi flujo de trabajo como
economista que:

\begin{itemize}
\tightlist
\item
  Mantiene múltiples blogs académicos (epsilon-y-beta, axiomata,
  methodica, etc.)
\item
  Desarrolla scripts de automatización en Python
\item
  Trabaja con análisis econométrico y ciencia de datos
\item
  Usa Quarto para publicaciones académicas
\item
  Necesita entornos reproducibles para colaboración
\end{itemize}

\section{Por Qué Miniconda en Lugar de
Anaconda}\label{por-quuxe9-miniconda-en-lugar-de-anaconda}

\subsection{Comparación}\label{comparaciuxf3n}

\begin{longtable}[]{@{}
  >{\raggedright\arraybackslash}p{(\linewidth - 4\tabcolsep) * \real{0.3265}}
  >{\raggedright\arraybackslash}p{(\linewidth - 4\tabcolsep) * \real{0.2245}}
  >{\raggedright\arraybackslash}p{(\linewidth - 4\tabcolsep) * \real{0.4490}}@{}}
\toprule\noalign{}
\begin{minipage}[b]{\linewidth}\raggedright
Característica
\end{minipage} & \begin{minipage}[b]{\linewidth}\raggedright
Miniconda
\end{minipage} & \begin{minipage}[b]{\linewidth}\raggedright
Anaconda Distribution
\end{minipage} \\
\midrule\noalign{}
\endhead
\bottomrule\noalign{}
\endlastfoot
Tamaño inicial & \textasciitilde50 MB & \textasciitilde3-5 GB \\
Paquetes preinstalados & Mínimos (conda, Python, pip) & 250+ paquetes \\
Control total & Sí & Limitado \\
Velocidad de instalación & Rápida & Lenta \\
Ideal para & Usuarios avanzados, servidores & Principiantes \\
\end{longtable}

\subsection{Ventajas}\label{ventajas}

\begin{enumerate}
\def\labelenumi{\arabic{enumi}.}
\tightlist
\item
  \textbf{Control Total:} Instalas solo lo que necesitas para cada
  proyecto
\item
  \textbf{Múltiples Entornos:} Ideal para tus 10+ blogs con dependencias
  diferentes
\item
  \textbf{Reproducibilidad:} Archivos \texttt{environment.yml} más
  pequeños y claros
\item
  \textbf{Velocidad:} Con Mamba, la resolución de dependencias es
  ultrarrápida
\item
  \textbf{Integración con Sistema:} Usa Quarto y R del sistema (pacman),
  Python de conda
\end{enumerate}

\section{Preparación del Sistema}\label{preparaciuxf3n-del-sistema}

\subsection{1. Verificar Instalaciones
Existentes}\label{verificar-instalaciones-existentes}

\begin{Shaded}
\begin{Highlighting}[]
\CommentTok{\# Verificar si conda ya está instalado}
\FunctionTok{which}\NormalTok{ conda}

\CommentTok{\# Verificar versión (si existe)}
\ExtensionTok{conda} \AttributeTok{{-}{-}version}

\CommentTok{\# Verificar entornos existentes}
\ExtensionTok{conda}\NormalTok{ env list}
\end{Highlighting}
\end{Shaded}

\subsection{2. Desinstalar Instalaciones Previas (Si Es
Necesario)}\label{desinstalar-instalaciones-previas-si-es-necesario}

\textbf{Si tienes Anaconda o Miniconda antiguo:}

\begin{Shaded}
\begin{Highlighting}[]
\CommentTok{\# Desactivar entorno base}
\ExtensionTok{conda}\NormalTok{ deactivate}

\CommentTok{\# Eliminar directorio de instalación}
\FunctionTok{rm} \AttributeTok{{-}rf}\NormalTok{ \textasciitilde{}/miniconda3}
\CommentTok{\# O si instalaste Anaconda:}
\FunctionTok{rm} \AttributeTok{{-}rf}\NormalTok{ \textasciitilde{}/anaconda3}

\CommentTok{\# Limpiar configuraciones (OPCIONAL {-} solo si quieres empezar de cero)}
\FunctionTok{rm} \AttributeTok{{-}rf}\NormalTok{ \textasciitilde{}/.conda}
\FunctionTok{rm} \AttributeTok{{-}f}\NormalTok{ \textasciitilde{}/.condarc}

\CommentTok{\# Limpiar inicialización de shells}
\CommentTok{\# Fish}
\FunctionTok{nano}\NormalTok{ \textasciitilde{}/.config/fish/config.fish  }\CommentTok{\# Eliminar bloque conda initialize}

\CommentTok{\# Bash}
\FunctionTok{nano}\NormalTok{ \textasciitilde{}/.bashrc  }\CommentTok{\# Eliminar bloque conda initialize}

\CommentTok{\# Zsh}
\FunctionTok{nano}\NormalTok{ \textasciitilde{}/.zshrc  }\CommentTok{\# Eliminar bloque conda initialize}
\end{Highlighting}
\end{Shaded}

\subsection{3. Verificar Herramientas del
Sistema}\label{verificar-herramientas-del-sistema}

\begin{Shaded}
\begin{Highlighting}[]
\CommentTok{\# Verificar Quarto (debe estar instalado)}
\ExtensionTok{quarto} \AttributeTok{{-}{-}version}

\CommentTok{\# Verificar R (debe estar instalado)}
\ExtensionTok{R} \AttributeTok{{-}{-}version}

\CommentTok{\# Verificar LaTeX (para PDFs en Quarto)}
\ExtensionTok{pdflatex} \AttributeTok{{-}{-}version}

\CommentTok{\# Si falta algo, instalar con pacman}
\FunctionTok{sudo}\NormalTok{ pacman }\AttributeTok{{-}S}\NormalTok{ quarto r texlive{-}most}
\end{Highlighting}
\end{Shaded}

\section{Instalación de Miniconda}\label{instalaciuxf3n-de-miniconda}

\subsection{Método 1: Instalación Rápida
(Recomendado)}\label{muxe9todo-1-instalaciuxf3n-ruxe1pida-recomendado}

\begin{Shaded}
\begin{Highlighting}[]
\CommentTok{\# 1. Crear directorio temporal}
\FunctionTok{mkdir} \AttributeTok{{-}p}\NormalTok{ \textasciitilde{}/miniconda3}

\CommentTok{\# 2. Descargar instalador (x86\_64 para Arch Linux estándar)}
\ExtensionTok{curl}\NormalTok{ https://repo.anaconda.com/miniconda/Miniconda3{-}latest{-}Linux{-}x86\_64.sh }\DataTypeTok{\textbackslash{}}
  \AttributeTok{{-}o}\NormalTok{ \textasciitilde{}/miniconda3/miniconda.sh}

\CommentTok{\# 3. Verificar integridad (RECOMENDADO)}
\FunctionTok{sha256sum}\NormalTok{ \textasciitilde{}/miniconda3/miniconda.sh}
\CommentTok{\# Comparar con el hash oficial en https://repo.anaconda.com/miniconda/}

\CommentTok{\# 4. Ejecutar instalador}
\FunctionTok{bash}\NormalTok{ \textasciitilde{}/miniconda3/miniconda.sh }\AttributeTok{{-}b} \AttributeTok{{-}u} \AttributeTok{{-}p}\NormalTok{ \textasciitilde{}/miniconda3}

\CommentTok{\# 5. Eliminar instalador}
\FunctionTok{rm}\NormalTok{ \textasciitilde{}/miniconda3/miniconda.sh}
\end{Highlighting}
\end{Shaded}

\textbf{opciones:}

\begin{itemize}
\tightlist
\item
  \texttt{-b}: Modo batch (sin prompts interactivos)
\item
  \texttt{-u}: Actualizar instalación existente si la hay
\item
  \texttt{-p\ \textasciitilde{}/miniconda3}: Ruta de instalación
\end{itemize}

\subsection{Método 2: Instalación
Interactiva}\label{muxe9todo-2-instalaciuxf3n-interactiva}

\begin{Shaded}
\begin{Highlighting}[]
\CommentTok{\# 1. Descargar instalador}
\FunctionTok{wget}\NormalTok{ https://repo.anaconda.com/miniconda/Miniconda3{-}latest{-}Linux{-}x86\_64.sh}

\CommentTok{\# 2. Ejecutar instalador interactivo}
\FunctionTok{bash}\NormalTok{ Miniconda3{-}latest{-}Linux{-}x86\_64.sh}

\CommentTok{\# Seguir instrucciones:}
\CommentTok{\# {-} Aceptar licencia: yes}
\CommentTok{\# {-} Ubicación: /home/achalmaedison/miniconda3 (default)}
\CommentTok{\# {-} Inicializar conda: yes (IMPORTANTE)}

\CommentTok{\# 3. Eliminar instalador}
\FunctionTok{rm}\NormalTok{ Miniconda3{-}latest{-}Linux{-}x86\_64.sh}
\end{Highlighting}
\end{Shaded}

\section{Configuración Inicial}\label{configuraciuxf3n-inicial}

\subsection{1. Inicializar Conda para Todos los
Shells}\label{inicializar-conda-para-todos-los-shells}

\begin{Shaded}
\begin{Highlighting}[]
\CommentTok{\# Activar conda temporalmente}
\BuiltInTok{source}\NormalTok{ \textasciitilde{}/miniconda3/bin/activate}

\CommentTok{\# Inicializar para todos los shells disponibles}
\ExtensionTok{conda}\NormalTok{ init }\AttributeTok{{-}{-}all}

\CommentTok{\# Resultado: Modifica automáticamente:}
\CommentTok{\# {-} \textasciitilde{}/.bashrc (Bash)}
\CommentTok{\# {-} \textasciitilde{}/.zshrc (Zsh)}
\CommentTok{\# {-} \textasciitilde{}/.config/fish/config.fish (Fish)}
\end{Highlighting}
\end{Shaded}

\subsection{2. Recargar Shell}\label{recargar-shell}

\begin{Shaded}
\begin{Highlighting}[]
\CommentTok{\# Fish}
\BuiltInTok{exec}\NormalTok{ fish}

\CommentTok{\# Bash}
\BuiltInTok{source}\NormalTok{ \textasciitilde{}/.bashrc}

\CommentTok{\# Zsh}
\BuiltInTok{source}\NormalTok{ \textasciitilde{}/.zshrc}
\end{Highlighting}
\end{Shaded}

\textbf{Verificación:} Deberías ver \texttt{(base)} al inicio de tu
prompt.

\subsection{3. Desactivar Auto-Activación de Base
(RECOMENDADO)}\label{desactivar-auto-activaciuxf3n-de-base-recomendado}

\begin{Shaded}
\begin{Highlighting}[]
\CommentTok{\# Evitar que base se active automáticamente}
\ExtensionTok{conda}\NormalTok{ config }\AttributeTok{{-}{-}set}\NormalTok{ auto\_activate\_base false}

\CommentTok{\# Recargar shell}
\BuiltInTok{exec}\NormalTok{ fish  }\CommentTok{\# o source \textasciitilde{}/.bashrc / source \textasciitilde{}/.zshrc}
\end{Highlighting}
\end{Shaded}

\textbf{Razón:} Es mejor práctica activar entornos específicos
manualmente.

\section{Instalación de Mamba}\label{instalaciuxf3n-de-mamba}

Mamba es un solver ultrarrápido para conda que acelera la instalación de
paquetes hasta 10x.

\subsection{1. Instalar Mamba en el Entorno
Base}\label{instalar-mamba-en-el-entorno-base}

\begin{Shaded}
\begin{Highlighting}[]
\CommentTok{\# Activar base}
\ExtensionTok{conda}\NormalTok{ activate base}

\CommentTok{\# Instalar mamba desde conda{-}forge}
\ExtensionTok{conda}\NormalTok{ install mamba }\AttributeTok{{-}c}\NormalTok{ conda{-}forge }\AttributeTok{{-}y}

\CommentTok{\# Configurar mamba como solver por defecto}
\ExtensionTok{conda}\NormalTok{ config }\AttributeTok{{-}{-}set}\NormalTok{ solver libmamba}
\end{Highlighting}
\end{Shaded}

\subsection{2. Verificar Instalación}\label{verificar-instalaciuxf3n}

\begin{Shaded}
\begin{Highlighting}[]
\CommentTok{\# Verificar versión}
\ExtensionTok{mamba} \AttributeTok{{-}{-}version}

\CommentTok{\# Debería mostrar algo como: mamba 1.5.x}
\end{Highlighting}
\end{Shaded}

\section{Configuración de Fish
Shell}\label{configuraciuxf3n-de-fish-shell}

\subsection{1. Inicialización
Automática}\label{inicializaciuxf3n-automuxe1tica}

\textbf{Conda ya debería haber configurado Fish automáticamente.}
Verificar:

\begin{Shaded}
\begin{Highlighting}[]
\NormalTok{\# Abrir configuración de Fish}
\NormalTok{nano \textasciitilde{}/.config/fish/config.fish}
\end{Highlighting}
\end{Shaded}

\textbf{Debe contener algo como:}

\begin{Shaded}
\begin{Highlighting}[]
\NormalTok{\# \textgreater{}\textgreater{}\textgreater{} conda initialize \textgreater{}\textgreater{}\textgreater{}}
\NormalTok{\# !! Contents within this block are managed by \textquotesingle{}conda init\textquotesingle{} !!}
\NormalTok{if test {-}f /home/achalmaedison/miniconda3/bin/conda}
\NormalTok{    eval /home/achalmaedison/miniconda3/bin/conda "shell.fish" "hook" $argv | source}
\NormalTok{else}
\NormalTok{    if test {-}f "/home/achalmaedison/miniconda3/etc/fish/conf.d/conda.fish"}
\NormalTok{        . "/home/achalmaedison/miniconda3/etc/fish/conf.d/conda.fish"}
\NormalTok{    else}
\NormalTok{        set {-}x PATH "/home/achalmaedison/miniconda3/bin" $PATH}
\NormalTok{    end}
\NormalTok{end}
\NormalTok{\# \textless{}\textless{}\textless{} conda initialize \textless{}\textless{}\textless{}}
\end{Highlighting}
\end{Shaded}

\subsection{2. Agregar Alias Útiles}\label{agregar-alias-uxfatiles}

\begin{Shaded}
\begin{Highlighting}[]
\NormalTok{\# Editar configuración de Fish}
\NormalTok{nano \textasciitilde{}/.config/fish/config.fish}
\end{Highlighting}
\end{Shaded}

\textbf{Agregar al final:}

\begin{Shaded}
\begin{Highlighting}[]
\NormalTok{\# ========================================================================}
\NormalTok{\# ALIASES DE CONDA Y PROYECTOS}
\NormalTok{\# ========================================================================}

\NormalTok{\# Gestión de entornos}
\NormalTok{alias ce="conda env list"}
\NormalTok{alias ca="conda activate"}
\NormalTok{alias cda="conda deactivate"}

\NormalTok{\# Blogs}
\NormalTok{alias blog="conda activate econblog \&\& cd \textasciitilde{}/Proyectos/epsilon{-}y{-}beta"}
\NormalTok{alias blogoff="conda deactivate"}

\NormalTok{\# Scripts}
\NormalTok{alias scripts="conda activate scripts{-}env \&\& cd \textasciitilde{}/Proyectos/scripts"}
\NormalTok{alias scriptsoff="conda deactivate"}

\NormalTok{\# Limpieza}
\NormalTok{alias conda{-}clean="conda clean {-}{-}all {-}y"}

\NormalTok{\# Actualización}
\NormalTok{alias conda{-}update="conda update {-}n base conda {-}y \&\& mamba update {-}{-}all {-}y"}
\end{Highlighting}
\end{Shaded}

\textbf{Recargar configuración:}

\begin{Shaded}
\begin{Highlighting}[]
\NormalTok{source \textasciitilde{}/.config/fish/config.fish}
\end{Highlighting}
\end{Shaded}

\subsection{3. Configuración para Bash/Zsh
(Opcional)}\label{configuraciuxf3n-para-bashzsh-opcional}

\textbf{Si también usas Bash o Zsh:}

\begin{Shaded}
\begin{Highlighting}[]
\CommentTok{\# Bash: \textasciitilde{}/.bashrc}
\CommentTok{\# Zsh: \textasciitilde{}/.zshrc}

\CommentTok{\# Agregar aliases}
\BuiltInTok{alias}\NormalTok{ ce=}\StringTok{"conda env list"}
\BuiltInTok{alias}\NormalTok{ ca=}\StringTok{"conda activate"}
\BuiltInTok{alias}\NormalTok{ cda=}\StringTok{"conda deactivate"}
\BuiltInTok{alias}\NormalTok{ blog=}\StringTok{"conda activate econblog \&\& cd \textasciitilde{}/Proyectos/epsilon{-}y{-}beta"}
\BuiltInTok{alias}\NormalTok{ scripts=}\StringTok{"conda activate scripts{-}env \&\& cd \textasciitilde{}/Proyectos/scripts"}
\BuiltInTok{alias}\NormalTok{ conda{-}clean=}\StringTok{"conda clean {-}{-}all {-}y"}
\BuiltInTok{alias}\NormalTok{ conda{-}update=}\StringTok{"conda update {-}n base conda {-}y \&\& mamba update {-}{-}all {-}y"}
\end{Highlighting}
\end{Shaded}

\section{Configuración Óptima de
Conda}\label{configuraciuxf3n-uxf3ptima-de-conda}

\subsection{1. Crear Archivo de
Configuración}\label{crear-archivo-de-configuraciuxf3n}

\begin{Shaded}
\begin{Highlighting}[]
\CommentTok{\# El archivo \textasciitilde{}/.condarc almacena la configuración global de conda}
\FunctionTok{nano}\NormalTok{ \textasciitilde{}/.condarc}
\end{Highlighting}
\end{Shaded}

\subsection{2. Mi Configuración}\label{mi-configuraciuxf3n}

\begin{Shaded}
\begin{Highlighting}[]
\CommentTok{\# ========================================================================}
\CommentTok{\# CONFIGURACIÓN GLOBAL DE CONDA}
\CommentTok{\# Usuario: Edison Achalma (achalmaedison)}
\CommentTok{\# Sistema: Arch Linux}
\CommentTok{\# Fecha: Enero 2026}
\CommentTok{\# ========================================================================}

\CommentTok{\# {-}{-}{-}{-}{-}{-}{-}{-}{-}{-}{-}{-}{-}{-}{-}{-}{-}{-}{-}{-}{-}{-}{-}{-}{-}{-}{-}{-}{-}{-}{-}{-}{-}{-}{-}{-}{-}{-}{-}{-}{-}{-}{-}{-}{-}{-}{-}{-}{-}{-}{-}{-}{-}{-}{-}{-}{-}{-}{-}{-}{-}{-}{-}{-}{-}{-}{-}{-}{-}{-}{-}{-}}
\CommentTok{\# Canales (Prioridad de búsqueda de paquetes)}
\CommentTok{\# {-}{-}{-}{-}{-}{-}{-}{-}{-}{-}{-}{-}{-}{-}{-}{-}{-}{-}{-}{-}{-}{-}{-}{-}{-}{-}{-}{-}{-}{-}{-}{-}{-}{-}{-}{-}{-}{-}{-}{-}{-}{-}{-}{-}{-}{-}{-}{-}{-}{-}{-}{-}{-}{-}{-}{-}{-}{-}{-}{-}{-}{-}{-}{-}{-}{-}{-}{-}{-}{-}{-}{-}}
\FunctionTok{channels}\KeywordTok{:}
\AttributeTok{  }\KeywordTok{{-}}\AttributeTok{ conda{-}forge}\CommentTok{        \# Canal principal (comunidad, actualizado)}
\AttributeTok{  }\KeywordTok{{-}}\AttributeTok{ defaults}\CommentTok{           \# Canal oficial de Anaconda}

\CommentTok{\# Prioridad estricta (usar en orden de la lista)}
\FunctionTok{channel\_priority}\KeywordTok{:}\AttributeTok{ strict}

\CommentTok{\# {-}{-}{-}{-}{-}{-}{-}{-}{-}{-}{-}{-}{-}{-}{-}{-}{-}{-}{-}{-}{-}{-}{-}{-}{-}{-}{-}{-}{-}{-}{-}{-}{-}{-}{-}{-}{-}{-}{-}{-}{-}{-}{-}{-}{-}{-}{-}{-}{-}{-}{-}{-}{-}{-}{-}{-}{-}{-}{-}{-}{-}{-}{-}{-}{-}{-}{-}{-}{-}{-}{-}{-}}
\CommentTok{\# Solver (Motor de resolución de dependencias)}
\CommentTok{\# {-}{-}{-}{-}{-}{-}{-}{-}{-}{-}{-}{-}{-}{-}{-}{-}{-}{-}{-}{-}{-}{-}{-}{-}{-}{-}{-}{-}{-}{-}{-}{-}{-}{-}{-}{-}{-}{-}{-}{-}{-}{-}{-}{-}{-}{-}{-}{-}{-}{-}{-}{-}{-}{-}{-}{-}{-}{-}{-}{-}{-}{-}{-}{-}{-}{-}{-}{-}{-}{-}{-}{-}}
\FunctionTok{solver}\KeywordTok{:}\AttributeTok{ libmamba}\CommentTok{      \# Usar mamba (mucho más rápido que classic)}

\CommentTok{\# {-}{-}{-}{-}{-}{-}{-}{-}{-}{-}{-}{-}{-}{-}{-}{-}{-}{-}{-}{-}{-}{-}{-}{-}{-}{-}{-}{-}{-}{-}{-}{-}{-}{-}{-}{-}{-}{-}{-}{-}{-}{-}{-}{-}{-}{-}{-}{-}{-}{-}{-}{-}{-}{-}{-}{-}{-}{-}{-}{-}{-}{-}{-}{-}{-}{-}{-}{-}{-}{-}{-}{-}}
\CommentTok{\# Activación Automática}
\CommentTok{\# {-}{-}{-}{-}{-}{-}{-}{-}{-}{-}{-}{-}{-}{-}{-}{-}{-}{-}{-}{-}{-}{-}{-}{-}{-}{-}{-}{-}{-}{-}{-}{-}{-}{-}{-}{-}{-}{-}{-}{-}{-}{-}{-}{-}{-}{-}{-}{-}{-}{-}{-}{-}{-}{-}{-}{-}{-}{-}{-}{-}{-}{-}{-}{-}{-}{-}{-}{-}{-}{-}{-}{-}}
\FunctionTok{auto\_activate\_base}\KeywordTok{:}\AttributeTok{ }\CharTok{false}\CommentTok{   \# NO activar base automáticamente}

\CommentTok{\# {-}{-}{-}{-}{-}{-}{-}{-}{-}{-}{-}{-}{-}{-}{-}{-}{-}{-}{-}{-}{-}{-}{-}{-}{-}{-}{-}{-}{-}{-}{-}{-}{-}{-}{-}{-}{-}{-}{-}{-}{-}{-}{-}{-}{-}{-}{-}{-}{-}{-}{-}{-}{-}{-}{-}{-}{-}{-}{-}{-}{-}{-}{-}{-}{-}{-}{-}{-}{-}{-}{-}{-}}
\CommentTok{\# Notificaciones y Actualizaciones}
\CommentTok{\# {-}{-}{-}{-}{-}{-}{-}{-}{-}{-}{-}{-}{-}{-}{-}{-}{-}{-}{-}{-}{-}{-}{-}{-}{-}{-}{-}{-}{-}{-}{-}{-}{-}{-}{-}{-}{-}{-}{-}{-}{-}{-}{-}{-}{-}{-}{-}{-}{-}{-}{-}{-}{-}{-}{-}{-}{-}{-}{-}{-}{-}{-}{-}{-}{-}{-}{-}{-}{-}{-}{-}{-}}
\FunctionTok{auto\_update\_conda}\KeywordTok{:}\AttributeTok{ }\CharTok{false}\CommentTok{    \# NO actualizar conda automáticamente}
\FunctionTok{notify\_outdated\_conda}\KeywordTok{:}\AttributeTok{ }\CharTok{true}\CommentTok{ \# Notificar si hay nueva versión}

\CommentTok{\# {-}{-}{-}{-}{-}{-}{-}{-}{-}{-}{-}{-}{-}{-}{-}{-}{-}{-}{-}{-}{-}{-}{-}{-}{-}{-}{-}{-}{-}{-}{-}{-}{-}{-}{-}{-}{-}{-}{-}{-}{-}{-}{-}{-}{-}{-}{-}{-}{-}{-}{-}{-}{-}{-}{-}{-}{-}{-}{-}{-}{-}{-}{-}{-}{-}{-}{-}{-}{-}{-}{-}{-}}
\CommentTok{\# Gestión de Paquetes}
\CommentTok{\# {-}{-}{-}{-}{-}{-}{-}{-}{-}{-}{-}{-}{-}{-}{-}{-}{-}{-}{-}{-}{-}{-}{-}{-}{-}{-}{-}{-}{-}{-}{-}{-}{-}{-}{-}{-}{-}{-}{-}{-}{-}{-}{-}{-}{-}{-}{-}{-}{-}{-}{-}{-}{-}{-}{-}{-}{-}{-}{-}{-}{-}{-}{-}{-}{-}{-}{-}{-}{-}{-}{-}{-}}
\FunctionTok{pip\_interop\_enabled}\KeywordTok{:}\AttributeTok{ }\CharTok{true}\CommentTok{   \# Permitir interoperabilidad con pip}
\FunctionTok{always\_yes}\KeywordTok{:}\AttributeTok{ }\CharTok{false}\CommentTok{           \# Pedir confirmación antes de instalar}

\CommentTok{\# {-}{-}{-}{-}{-}{-}{-}{-}{-}{-}{-}{-}{-}{-}{-}{-}{-}{-}{-}{-}{-}{-}{-}{-}{-}{-}{-}{-}{-}{-}{-}{-}{-}{-}{-}{-}{-}{-}{-}{-}{-}{-}{-}{-}{-}{-}{-}{-}{-}{-}{-}{-}{-}{-}{-}{-}{-}{-}{-}{-}{-}{-}{-}{-}{-}{-}{-}{-}{-}{-}{-}{-}}
\CommentTok{\# Almacenamiento}
\CommentTok{\# {-}{-}{-}{-}{-}{-}{-}{-}{-}{-}{-}{-}{-}{-}{-}{-}{-}{-}{-}{-}{-}{-}{-}{-}{-}{-}{-}{-}{-}{-}{-}{-}{-}{-}{-}{-}{-}{-}{-}{-}{-}{-}{-}{-}{-}{-}{-}{-}{-}{-}{-}{-}{-}{-}{-}{-}{-}{-}{-}{-}{-}{-}{-}{-}{-}{-}{-}{-}{-}{-}{-}{-}}
\FunctionTok{pkgs\_dirs}\KeywordTok{:}
\AttributeTok{  }\KeywordTok{{-}}\AttributeTok{ \textasciitilde{}/.conda/pkgs}\CommentTok{           \# Directorio de caché de paquetes}

\FunctionTok{envs\_dirs}\KeywordTok{:}
\AttributeTok{  }\KeywordTok{{-}}\AttributeTok{ \textasciitilde{}/miniconda3/envs}\CommentTok{       \# Directorio de entornos}
\AttributeTok{  }\KeywordTok{{-}}\AttributeTok{ \textasciitilde{}/.conda/envs}\CommentTok{           \# Directorio adicional de entornos}

\CommentTok{\# {-}{-}{-}{-}{-}{-}{-}{-}{-}{-}{-}{-}{-}{-}{-}{-}{-}{-}{-}{-}{-}{-}{-}{-}{-}{-}{-}{-}{-}{-}{-}{-}{-}{-}{-}{-}{-}{-}{-}{-}{-}{-}{-}{-}{-}{-}{-}{-}{-}{-}{-}{-}{-}{-}{-}{-}{-}{-}{-}{-}{-}{-}{-}{-}{-}{-}{-}{-}{-}{-}{-}{-}}
\CommentTok{\# Interfaz de Usuario}
\CommentTok{\# {-}{-}{-}{-}{-}{-}{-}{-}{-}{-}{-}{-}{-}{-}{-}{-}{-}{-}{-}{-}{-}{-}{-}{-}{-}{-}{-}{-}{-}{-}{-}{-}{-}{-}{-}{-}{-}{-}{-}{-}{-}{-}{-}{-}{-}{-}{-}{-}{-}{-}{-}{-}{-}{-}{-}{-}{-}{-}{-}{-}{-}{-}{-}{-}{-}{-}{-}{-}{-}{-}{-}{-}}
\FunctionTok{show\_channel\_urls}\KeywordTok{:}\AttributeTok{ }\CharTok{true}\CommentTok{     \# Mostrar URL del canal en listados}
\FunctionTok{json}\KeywordTok{:}\AttributeTok{ }\CharTok{false}\CommentTok{                 \# Salida en texto (no JSON)}
\FunctionTok{quiet}\KeywordTok{:}\AttributeTok{ }\CharTok{false}\CommentTok{                \# Mostrar mensajes normales}

\CommentTok{\# {-}{-}{-}{-}{-}{-}{-}{-}{-}{-}{-}{-}{-}{-}{-}{-}{-}{-}{-}{-}{-}{-}{-}{-}{-}{-}{-}{-}{-}{-}{-}{-}{-}{-}{-}{-}{-}{-}{-}{-}{-}{-}{-}{-}{-}{-}{-}{-}{-}{-}{-}{-}{-}{-}{-}{-}{-}{-}{-}{-}{-}{-}{-}{-}{-}{-}{-}{-}{-}{-}{-}{-}}
\CommentTok{\# Seguridad}
\CommentTok{\# {-}{-}{-}{-}{-}{-}{-}{-}{-}{-}{-}{-}{-}{-}{-}{-}{-}{-}{-}{-}{-}{-}{-}{-}{-}{-}{-}{-}{-}{-}{-}{-}{-}{-}{-}{-}{-}{-}{-}{-}{-}{-}{-}{-}{-}{-}{-}{-}{-}{-}{-}{-}{-}{-}{-}{-}{-}{-}{-}{-}{-}{-}{-}{-}{-}{-}{-}{-}{-}{-}{-}{-}}
\FunctionTok{ssl\_verify}\KeywordTok{:}\AttributeTok{ }\CharTok{true}\CommentTok{            \# Verificar certificados SSL}

\CommentTok{\# {-}{-}{-}{-}{-}{-}{-}{-}{-}{-}{-}{-}{-}{-}{-}{-}{-}{-}{-}{-}{-}{-}{-}{-}{-}{-}{-}{-}{-}{-}{-}{-}{-}{-}{-}{-}{-}{-}{-}{-}{-}{-}{-}{-}{-}{-}{-}{-}{-}{-}{-}{-}{-}{-}{-}{-}{-}{-}{-}{-}{-}{-}{-}{-}{-}{-}{-}{-}{-}{-}{-}{-}}
\CommentTok{\# Proxy (Descomentar si usas proxy corporativo)}
\CommentTok{\# {-}{-}{-}{-}{-}{-}{-}{-}{-}{-}{-}{-}{-}{-}{-}{-}{-}{-}{-}{-}{-}{-}{-}{-}{-}{-}{-}{-}{-}{-}{-}{-}{-}{-}{-}{-}{-}{-}{-}{-}{-}{-}{-}{-}{-}{-}{-}{-}{-}{-}{-}{-}{-}{-}{-}{-}{-}{-}{-}{-}{-}{-}{-}{-}{-}{-}{-}{-}{-}{-}{-}{-}}
\CommentTok{\# proxy\_servers:}
\CommentTok{\#   http: http://proxy.empresa.com:8080}
\CommentTok{\#   https: https://proxy.empresa.com:8080}
\end{Highlighting}
\end{Shaded}

\subsection{3. Verificar
Configuración}\label{verificar-configuraciuxf3n}

\begin{Shaded}
\begin{Highlighting}[]
\CommentTok{\# Ver configuración actual}
\ExtensionTok{conda}\NormalTok{ config }\AttributeTok{{-}{-}show}

\CommentTok{\# Ver solo configuraciones modificadas}
\ExtensionTok{conda}\NormalTok{ config }\AttributeTok{{-}{-}show{-}sources}
\end{Highlighting}
\end{Shaded}

\section{Integración con Herramientas del
Sistema}\label{integraciuxf3n-con-herramientas-del-sistema}

\subsection{1. Usar Quarto del Sistema (No de
Conda)}\label{usar-quarto-del-sistema-no-de-conda}

\textbf{Ventaja:} Quarto de pacman es más estable y actualizado.

\begin{Shaded}
\begin{Highlighting}[]
\CommentTok{\# Verificar que Quarto está en PATH del sistema}
\FunctionTok{which}\NormalTok{ quarto}
\CommentTok{\# Debe mostrar: /usr/bin/quarto}

\CommentTok{\# Verificar versión}
\ExtensionTok{quarto} \AttributeTok{{-}{-}version}
\end{Highlighting}
\end{Shaded}

\textbf{En entornos conda, Quarto del sistema se usará automáticamente.}

\subsection{2. Usar R del Sistema (No de
Conda)}\label{usar-r-del-sistema-no-de-conda}

\textbf{Ventaja:} RStudio funciona mejor con R de pacman.

\begin{Shaded}
\begin{Highlighting}[]
\CommentTok{\# Verificar R del sistema}
\FunctionTok{which}\NormalTok{ R}
\CommentTok{\# Debe mostrar: /usr/bin/R}

\CommentTok{\# Verificar versión}
\ExtensionTok{R} \AttributeTok{{-}{-}version}
\end{Highlighting}
\end{Shaded}

\textbf{R del sistema estará disponible en todos los entornos conda.}

\subsection{3. Usar LaTeX del Sistema}\label{usar-latex-del-sistema}

\begin{Shaded}
\begin{Highlighting}[]
\CommentTok{\# Verificar LaTeX}
\FunctionTok{which}\NormalTok{ pdflatex}
\CommentTok{\# Debe mostrar: /usr/bin/pdflatex}

\CommentTok{\# Verificar distribución}
\ExtensionTok{pdflatex} \AttributeTok{{-}{-}version}
\end{Highlighting}
\end{Shaded}

\section{Verificación de la
Instalación}\label{verificaciuxf3n-de-la-instalaciuxf3n}

\subsection{1. Verificación Básica}\label{verificaciuxf3n-buxe1sica}

\begin{Shaded}
\begin{Highlighting}[]
\CommentTok{\# Versión de conda}
\ExtensionTok{conda} \AttributeTok{{-}{-}version}

\CommentTok{\# Versión de mamba}
\ExtensionTok{mamba} \AttributeTok{{-}{-}version}

\CommentTok{\# Versión de Python (en base)}
\ExtensionTok{python} \AttributeTok{{-}{-}version}

\CommentTok{\# Información del sistema}
\ExtensionTok{conda}\NormalTok{ info}
\end{Highlighting}
\end{Shaded}

\subsection{2. Verificación de
Configuración}\label{verificaciuxf3n-de-configuraciuxf3n}

\begin{Shaded}
\begin{Highlighting}[]
\CommentTok{\# Ver configuración}
\ExtensionTok{conda}\NormalTok{ config }\AttributeTok{{-}{-}show}

\CommentTok{\# Ver canales}
\ExtensionTok{conda}\NormalTok{ config }\AttributeTok{{-}{-}show}\NormalTok{ channels}

\CommentTok{\# Ver solver}
\ExtensionTok{conda}\NormalTok{ config }\AttributeTok{{-}{-}show}\NormalTok{ solver}
\end{Highlighting}
\end{Shaded}

\subsection{3. Prueba de Funcionamiento}\label{prueba-de-funcionamiento}

\begin{Shaded}
\begin{Highlighting}[]
\CommentTok{\# Crear entorno de prueba}
\ExtensionTok{mamba}\NormalTok{ create }\AttributeTok{{-}n}\NormalTok{ test{-}env python=3.12 pandas }\AttributeTok{{-}y}

\CommentTok{\# Activar entorno}
\ExtensionTok{conda}\NormalTok{ activate test{-}env}

\CommentTok{\# Verificar instalación}
\ExtensionTok{python} \AttributeTok{{-}c} \StringTok{"import pandas; print(f\textquotesingle{}Pandas \{pandas.\_\_version\_\_\} funcionando correctamente\textquotesingle{})"}

\CommentTok{\# Desactivar}
\ExtensionTok{conda}\NormalTok{ deactivate}

\CommentTok{\# Eliminar entorno de prueba}
\ExtensionTok{conda}\NormalTok{ env remove }\AttributeTok{{-}n}\NormalTok{ test{-}env }\AttributeTok{{-}y}
\end{Highlighting}
\end{Shaded}

\begin{center}\rule{0.5\linewidth}{0.5pt}\end{center}

\section{Resolución de Problemas}\label{resoluciuxf3n-de-problemas}

\subsection{Problema 1: ``conda: command not
found''}\label{problema-1-conda-command-not-found}

\textbf{Causa:} Conda no está inicializado en tu shell.

\textbf{Solución:}

\begin{Shaded}
\begin{Highlighting}[]
\CommentTok{\# Inicializar manualmente}
\BuiltInTok{source}\NormalTok{ \textasciitilde{}/miniconda3/bin/activate}
\ExtensionTok{conda}\NormalTok{ init fish  }\CommentTok{\# o bash / zsh}

\CommentTok{\# Recargar shell}
\BuiltInTok{exec}\NormalTok{ fish}
\end{Highlighting}
\end{Shaded}

\subsection{Problema 2: ``(base) No Aparece en el
Prompt''}\label{problema-2-base-no-aparece-en-el-prompt}

\textbf{Causa:} \texttt{auto\_activate\_base} está en \texttt{false}.

\textbf{Solución (si quieres que aparezca):}

\begin{Shaded}
\begin{Highlighting}[]
\ExtensionTok{conda}\NormalTok{ config }\AttributeTok{{-}{-}set}\NormalTok{ auto\_activate\_base true}
\BuiltInTok{exec}\NormalTok{ fish}
\end{Highlighting}
\end{Shaded}

\textbf{O activar manualmente:}

\begin{Shaded}
\begin{Highlighting}[]
\ExtensionTok{conda}\NormalTok{ activate base}
\end{Highlighting}
\end{Shaded}

\subsection{Problema 3: ``Mamba No
Funciona''}\label{problema-3-mamba-no-funciona}

\textbf{Causa:} Mamba no está instalado o no está en PATH.

\textbf{Solución:}

\begin{Shaded}
\begin{Highlighting}[]
\CommentTok{\# Verificar instalación}
\ExtensionTok{conda}\NormalTok{ activate base}
\ExtensionTok{conda}\NormalTok{ list }\KeywordTok{|} \FunctionTok{grep}\NormalTok{ mamba}

\CommentTok{\# Si no está, instalar}
\ExtensionTok{conda}\NormalTok{ install mamba }\AttributeTok{{-}c}\NormalTok{ conda{-}forge }\AttributeTok{{-}y}

\CommentTok{\# Configurar como solver}
\ExtensionTok{conda}\NormalTok{ config }\AttributeTok{{-}{-}set}\NormalTok{ solver libmamba}
\end{Highlighting}
\end{Shaded}

\subsection{Problema 4: ``Instalación de Paquetes es Muy
Lenta''}\label{problema-4-instalaciuxf3n-de-paquetes-es-muy-lenta}

\textbf{Causa:} Usando solver clásico en lugar de mamba.

\textbf{Solución:}

\begin{Shaded}
\begin{Highlighting}[]
\CommentTok{\# Configurar mamba como solver}
\ExtensionTok{conda}\NormalTok{ config }\AttributeTok{{-}{-}set}\NormalTok{ solver libmamba}

\CommentTok{\# O usar mamba directamente}
\ExtensionTok{mamba}\NormalTok{ install paquete{-}nombre}
\end{Highlighting}
\end{Shaded}

\subsection{Problema 5: ``Conflictos de
Dependencias''}\label{problema-5-conflictos-de-dependencias}

\textbf{Causa:} Paquetes incompatibles en el mismo entorno.

\textbf{Solución:}

\begin{Shaded}
\begin{Highlighting}[]
\CommentTok{\# Opción 1: Usar mamba (mejor resolución)}
\ExtensionTok{mamba}\NormalTok{ install paquete{-}nombre}

\CommentTok{\# Opción 2: Crear entorno nuevo limpio}
\ExtensionTok{mamba}\NormalTok{ create }\AttributeTok{{-}n}\NormalTok{ nuevo{-}env python=3.12 paquete{-}nombre}

\CommentTok{\# Opción 3: Actualizar todos los paquetes}
\ExtensionTok{mamba}\NormalTok{ update }\AttributeTok{{-}{-}all}
\end{Highlighting}
\end{Shaded}

\subsection{Problema 6: ``Espacio en Disco
Lleno''}\label{problema-6-espacio-en-disco-lleno}

\textbf{Causa:} Caché de paquetes acumulado.

\textbf{Solución:}

\begin{Shaded}
\begin{Highlighting}[]
\CommentTok{\# Limpiar caché de paquetes}
\ExtensionTok{conda}\NormalTok{ clean }\AttributeTok{{-}{-}all} \AttributeTok{{-}y}

\CommentTok{\# Ver espacio usado}
\FunctionTok{du} \AttributeTok{{-}sh}\NormalTok{ \textasciitilde{}/.conda/pkgs}

\CommentTok{\# Eliminar entornos no usados}
\ExtensionTok{conda}\NormalTok{ env list}
\ExtensionTok{conda}\NormalTok{ env remove }\AttributeTok{{-}n}\NormalTok{ nombre{-}entorno{-}viejo}
\end{Highlighting}
\end{Shaded}

\subsection{Problema 7: ``Fish No Reconoce
Conda''}\label{problema-7-fish-no-reconoce-conda}

\textbf{Causa:} Inicialización incorrecta de Fish.

\textbf{Solución:}

\begin{Shaded}
\begin{Highlighting}[]
\CommentTok{\# Verificar bloque de inicialización}
\FunctionTok{cat}\NormalTok{ \textasciitilde{}/.config/fish/config.fish }\KeywordTok{|} \FunctionTok{grep}\NormalTok{ conda}

\CommentTok{\# Si no existe, agregar manualmente}
\FunctionTok{nano}\NormalTok{ \textasciitilde{}/.config/fish/config.fish}
\end{Highlighting}
\end{Shaded}

\textbf{Agregar:}

\begin{Shaded}
\begin{Highlighting}[]
\NormalTok{\# \textgreater{}\textgreater{}\textgreater{} conda initialize \textgreater{}\textgreater{}\textgreater{}}
\NormalTok{eval /home/achalmaedison/miniconda3/bin/conda "shell.fish" "hook" $argv | source}
\NormalTok{\# \textless{}\textless{}\textless{} conda initialize \textless{}\textless{}\textless{}}
\end{Highlighting}
\end{Shaded}

\textbf{Recargar:}

\begin{Shaded}
\begin{Highlighting}[]
\NormalTok{exec fish}
\end{Highlighting}
\end{Shaded}

\begin{center}\rule{0.5\linewidth}{0.5pt}\end{center}

\section{Comandos de Referencia
Rápida}\label{comandos-de-referencia-ruxe1pida}

\subsection{Gestión de Conda}\label{gestiuxf3n-de-conda}

\begin{Shaded}
\begin{Highlighting}[]
\CommentTok{\# Actualizar conda}
\ExtensionTok{conda}\NormalTok{ update }\AttributeTok{{-}n}\NormalTok{ base conda }\AttributeTok{{-}y}

\CommentTok{\# Actualizar mamba}
\ExtensionTok{mamba}\NormalTok{ update mamba }\AttributeTok{{-}y}

\CommentTok{\# Información del sistema}
\ExtensionTok{conda}\NormalTok{ info}

\CommentTok{\# Ver configuración}
\ExtensionTok{conda}\NormalTok{ config }\AttributeTok{{-}{-}show}

\CommentTok{\# Limpiar caché}
\ExtensionTok{conda}\NormalTok{ clean }\AttributeTok{{-}{-}all} \AttributeTok{{-}y}
\end{Highlighting}
\end{Shaded}

\subsection{Gestión de Entornos}\label{gestiuxf3n-de-entornos}

\begin{Shaded}
\begin{Highlighting}[]
\CommentTok{\# Listar entornos}
\ExtensionTok{conda}\NormalTok{ env list}

\CommentTok{\# Crear entorno}
\ExtensionTok{mamba}\NormalTok{ create }\AttributeTok{{-}n}\NormalTok{ nombre{-}env python=3.12 paquete1 paquete2 }\AttributeTok{{-}y}

\CommentTok{\# Activar entorno}
\ExtensionTok{conda}\NormalTok{ activate nombre{-}env}

\CommentTok{\# Desactivar entorno}
\ExtensionTok{conda}\NormalTok{ deactivate}

\CommentTok{\# Eliminar entorno}
\ExtensionTok{conda}\NormalTok{ env remove }\AttributeTok{{-}n}\NormalTok{ nombre{-}env }\AttributeTok{{-}y}

\CommentTok{\# Exportar entorno}
\ExtensionTok{conda}\NormalTok{ env export }\AttributeTok{{-}{-}no{-}builds} \OperatorTok{\textgreater{}}\NormalTok{ environment.yml}

\CommentTok{\# Crear desde archivo}
\ExtensionTok{conda}\NormalTok{ env create }\AttributeTok{{-}f}\NormalTok{ environment.yml}
\end{Highlighting}
\end{Shaded}

\subsection{Gestión de Paquetes}\label{gestiuxf3n-de-paquetes}

\begin{Shaded}
\begin{Highlighting}[]
\CommentTok{\# Listar paquetes (entorno activo)}
\ExtensionTok{conda}\NormalTok{ list}

\CommentTok{\# Buscar paquete}
\ExtensionTok{conda}\NormalTok{ search nombre{-}paquete}

\CommentTok{\# Instalar paquete}
\ExtensionTok{mamba}\NormalTok{ install nombre{-}paquete }\AttributeTok{{-}y}

\CommentTok{\# Instalar desde canal específico}
\ExtensionTok{mamba}\NormalTok{ install }\AttributeTok{{-}c}\NormalTok{ conda{-}forge nombre{-}paquete }\AttributeTok{{-}y}

\CommentTok{\# Instalar versión específica}
\ExtensionTok{mamba}\NormalTok{ install nombre{-}paquete=1.2.3 }\AttributeTok{{-}y}

\CommentTok{\# Actualizar paquete}
\ExtensionTok{mamba}\NormalTok{ update nombre{-}paquete }\AttributeTok{{-}y}

\CommentTok{\# Actualizar todos los paquetes}
\ExtensionTok{mamba}\NormalTok{ update }\AttributeTok{{-}{-}all} \AttributeTok{{-}y}

\CommentTok{\# Desinstalar paquete}
\ExtensionTok{conda}\NormalTok{ remove nombre{-}paquete }\AttributeTok{{-}y}
\end{Highlighting}
\end{Shaded}

\subsection{Uso con Pip}\label{uso-con-pip}

\begin{Shaded}
\begin{Highlighting}[]
\CommentTok{\# Instalar con pip (dentro de entorno conda)}
\ExtensionTok{pip}\NormalTok{ install nombre{-}paquete}

\CommentTok{\# Listar paquetes pip}
\ExtensionTok{pip}\NormalTok{ list}

\CommentTok{\# Actualizar paquete pip}
\ExtensionTok{pip}\NormalTok{ install }\AttributeTok{{-}{-}upgrade}\NormalTok{ nombre{-}paquete}

\CommentTok{\# Desinstalar paquete pip}
\ExtensionTok{pip}\NormalTok{ uninstall nombre{-}paquete}
\end{Highlighting}
\end{Shaded}

\subsection{Información y Debugging}\label{informaciuxf3n-y-debugging}

\begin{Shaded}
\begin{Highlighting}[]
\CommentTok{\# Información de entorno activo}
\ExtensionTok{conda}\NormalTok{ info }\AttributeTok{{-}{-}envs}

\CommentTok{\# Información de paquete específico}
\ExtensionTok{conda}\NormalTok{ list nombre{-}paquete}

\CommentTok{\# Ver dependencias de paquete}
\ExtensionTok{conda}\NormalTok{ search nombre{-}paquete }\AttributeTok{{-}{-}info}

\CommentTok{\# Verificar integridad de entorno}
\ExtensionTok{conda}\NormalTok{ doctor}

\CommentTok{\# Modo verbose (más información)}
\ExtensionTok{conda}\NormalTok{ install }\AttributeTok{{-}{-}debug}\NormalTok{ nombre{-}paquete}
\end{Highlighting}
\end{Shaded}

\section{Recursos Adicionales}\label{recursos-adicionales}

\textbf{Documentación Oficial:}

\begin{itemize}
\tightlist
\item
  Conda: https://docs.conda.io/
\item
  Mamba: https://mamba.readthedocs.io/
\item
  Conda-Forge: https://conda-forge.org/
\end{itemize}

\textbf{Comunidad:}

\begin{itemize}
\tightlist
\item
  Conda Discourse: https://conda.discourse.group/
\item
  GitHub Issues: https://github.com/conda/conda/issues
\end{itemize}

\section{Publicaciones Similares}\label{publicaciones-similares}

Si te interesó este artículo, te recomendamos que explores otros blogs y
recursos relacionados que pueden ampliar tus conocimientos. Aquí te dejo
algunas sugerencias:

\begin{enumerate}
\def\labelenumi{\arabic{enumi}.}
\tightlist
\item
  \href{https://numerus-scriptum.netlify.app/python/2020-06-19-instalacion-de-anaconda/index.pdf}{\faIcon{file-pdf}}
  \href{https://numerus-scriptum.netlify.app/python/2020-06-19-instalacion-de-anaconda}{Instalacion
  De Anaconda}
\item
  \href{https://numerus-scriptum.netlify.app/python/2020-06-20-configurar-entorno-virtual-python-anaconda/index.pdf}{\faIcon{file-pdf}}
  \href{https://numerus-scriptum.netlify.app/python/2020-06-20-configurar-entorno-virtual-python-anaconda}{Configurar
  Entorno Virtual Python Anaconda}
\item
  \href{https://numerus-scriptum.netlify.app/python/2021-04-17-01-introducion-a-la-programacion-con-python/index.pdf}{\faIcon{file-pdf}}
  \href{https://numerus-scriptum.netlify.app/python/2021-04-17-01-introducion-a-la-programacion-con-python}{01
  Introducion A La Programacion Con Python}
\item
  \href{https://numerus-scriptum.netlify.app/python/2021-05-31-02-variables-expresiones-y-statements-con-python/index.pdf}{\faIcon{file-pdf}}
  \href{https://numerus-scriptum.netlify.app/python/2021-05-31-02-variables-expresiones-y-statements-con-python}{02
  Variables Expresiones Y Statements Con Python}
\item
  \href{https://numerus-scriptum.netlify.app/python/2021-06-07-03-objetos-de-python/index.pdf}{\faIcon{file-pdf}}
  \href{https://numerus-scriptum.netlify.app/python/2021-06-07-03-objetos-de-python}{03
  Objetos De Python}
\item
  \href{https://numerus-scriptum.netlify.app/python/2021-06-14-04-ejecucion-condicional-con-python/index.pdf}{\faIcon{file-pdf}}
  \href{https://numerus-scriptum.netlify.app/python/2021-06-14-04-ejecucion-condicional-con-python}{04
  Ejecucion Condicional Con Python}
\item
  \href{https://numerus-scriptum.netlify.app/python/2021-06-21-05-iteraciones-con-python/index.pdf}{\faIcon{file-pdf}}
  \href{https://numerus-scriptum.netlify.app/python/2021-06-21-05-iteraciones-con-python}{05
  Iteraciones Con Python}
\item
  \href{https://numerus-scriptum.netlify.app/python/2021-08-16-06-funciones-con-python/index.pdf}{\faIcon{file-pdf}}
  \href{https://numerus-scriptum.netlify.app/python/2021-08-16-06-funciones-con-python}{06
  Funciones Con Python}
\item
  \href{https://numerus-scriptum.netlify.app/python/2021-08-23-07-dataframes-con-python/index.pdf}{\faIcon{file-pdf}}
  \href{https://numerus-scriptum.netlify.app/python/2021-08-23-07-dataframes-con-python}{07
  Dataframes Con Python}
\item
  \href{https://numerus-scriptum.netlify.app/python/2021-11-29-08-prediccion-y-metrica-de-performance-con-python/index.pdf}{\faIcon{file-pdf}}
  \href{https://numerus-scriptum.netlify.app/python/2021-11-29-08-prediccion-y-metrica-de-performance-con-python}{08
  Prediccion Y Metrica De Performance Con Python}
\item
  \href{https://numerus-scriptum.netlify.app/python/2021-12-06-09-metodos-de-machine-learning-para-clasificacion-con-python/index.pdf}{\faIcon{file-pdf}}
  \href{https://numerus-scriptum.netlify.app/python/2021-12-06-09-metodos-de-machine-learning-para-clasificacion-con-python}{09
  Metodos De Machine Learning Para Clasificacion Con Python}
\item
  \href{https://numerus-scriptum.netlify.app/python/2021-12-13-10-metodos-de-machine-learning-para-regresion-con-python/index.pdf}{\faIcon{file-pdf}}
  \href{https://numerus-scriptum.netlify.app/python/2021-12-13-10-metodos-de-machine-learning-para-regresion-con-python}{10
  Metodos De Machine Learning Para Regresion Con Python}
\item
  \href{https://numerus-scriptum.netlify.app/python/2022-10-31-11-validacion-cruzada-y-composicion-del-modelo-con-python/index.pdf}{\faIcon{file-pdf}}
  \href{https://numerus-scriptum.netlify.app/python/2022-10-31-11-validacion-cruzada-y-composicion-del-modelo-con-python}{11
  Validacion Cruzada Y Composicion Del Modelo Con Python}
\item
  \href{https://numerus-scriptum.netlify.app/python/2025-05-10-visualizacion-de-datos-con-python/index.pdf}{\faIcon{file-pdf}}
  \href{https://numerus-scriptum.netlify.app/python/2025-05-10-visualizacion-de-datos-con-python}{Visualizacion
  De Datos Con Python}
\end{enumerate}

Esperamos que encuentres estas publicaciones igualmente interesantes y
útiles. ¡Disfruta de la lectura!






\end{document}
